%% ma: L10-B15-0104
%% lop: 10
%% chuong: 6
%% bai: 15
%% dang: 10.15.4
%% loai: TN4
%% muc: VD
%% dap_an: "D"
%% nguon: "Ảnh giáo viên gửi 25/11/2026 (lô 2): Bài tập trắc nghiệm – Toán 10 Bài 15 Hàm số (sự biến thiên, đồ thị), Phần I Câu 15"
%% kien_thuc: "Lớp 10 Bài 15 (sự đồng biến, nghịch biến của hàm số)"
%% trang_thai: chinh-thuc
\begin{ex}
Xét tính đồng biến, nghịch biến của hàm số $f(x)=\dfrac{x-3}{x+5}$ trên khoảng $(-\infty;-5)$ và trên khoảng $(-5;+\infty)$. Khẳng định nào sau đây đúng?
\choice{Hàm số nghịch biến trên $(-\infty;-5)$, đồng biến trên $(-5;+\infty)$.}{Hàm số đồng biến trên $(-\infty;-5)$, nghịch biến trên $(-5;+\infty)$.}{Hàm số nghịch biến trên các khoảng $(-\infty;-5)$ và $(-5;+\infty)$.}{\True Hàm số đồng biến trên các khoảng $(-\infty;-5)$ và $(-5;+\infty)$.}
\loigiai{Với $x_1\ne x_2$ cùng thuộc một khoảng: $\dfrac{f(x_1)-f(x_2)}{x_1-x_2}=\dfrac{8}{(x_1+5)(x_2+5)}>0$ vì $x_1+5$, $x_2+5$ cùng dấu. Vậy hàm số đồng biến trên từng khoảng.}
\end{ex}
